\documentclass[UTF8,a4paper,11pt]{ctexart}
\usepackage{amsmath,amssymb}
\usepackage[margin=2.5cm]{geometry}
\usepackage[hidelinks]{hyperref}
\title{2025年全国Ⅱ卷第18题：唯一极值、零点与对称差}
\date{2026年10月24日}
\begin{document}
\maketitle
\noindent\textbf{固定来源：}\url{https://github.com/DxAThing/Gaokao-Math-Problems-Compilation/blob/f1ae169b58a75b3db6e35c670b3f0d5c78718188/content/2025/national_paper_2.tex#L496}

设 $0<k<1/3$，
\[
f(x)=\ln(1+x)-x+\frac{x^2}{2}-kx^3,\quad x>-1.
\]
仅在 $x>0$ 上讨论极值点 $x_1$ 和正零点 $x_2$。
\paragraph{一、唯一性}
\[
f'(x)=x^2\left(\frac1{1+x}-3k\right).
\]
令
\[
c=\frac1{3k}-1>0.
\]
对 $0<x<c$，$f'(x)>0$；对 $x>c$，$f'(x)<0$。
所以 $x_1=c$ 为唯一极值点，且是严格极大值点。
又 $f(0)=0$、$f(c)>0$，且
$\lim_{x\to\infty}f(x)=-\infty$。
由连续性及递减性，在 $(c,+\infty)$ 上有且仅有一个零点
$x_2$。注意 $x=0$ 不在题目要求的 $(0,+\infty)$ 内。

\paragraph{二、对称差严格递减}
令
\[
g(t)=f(c+t)-f(c-t),\qquad 0<t<c.
\]
链式法则给出正确的\textbf{加号}：
\[
g'(t)=f'(c+t)+f'(c-t).
\]
利用 $3k=1/(1+c)$，得
\[
f'(x)=\frac{x^2(c-x)}{(1+x)(1+c)}.
\]
令 $H(x)=x^2/(1+x)$，则
\[
H'(x)=\frac{x(x+2)}{(1+x)^2}>0\quad(x>0).
\]
因而
\[
g'(t)=\frac{t}{1+c}\bigl[H(c-t)-H(c+t)\bigr]<0,
\]
故 $g$ 在 $(0,c)$ 严格递减。

\paragraph{三、比较两零点位置}
由连续延拓的 $g(0)=0$ 及严格递减性，
\[
g(c)=f(2c)-f(0)=f(2c)<0.
\]
结合 $f(c)>0$，及 $f$ 在 $(c,\infty)$ 严格递减，
唯一正零点 $x_2$ 满足
\[
\boxed{x_1<x_2<2x_1}.
\]

\bigskip
\noindent\textbf{源码质检提示：}
上游解答中一行把 $g'(t)$ 误写为两个 $f'$ 之差；
链式法则正确结果是上文的\textbf{和}。
本稿为独立数学证明，尚未实际 XeLaTeX 编译。
\end{document}
